\subsection{The common-orientation quotient of the marked datum}
\label{subsec:v28-common-orientation}

There are two different operations on signed observations: changing one
orientation convention for the entire experiment, and discarding signs from
individual observations.  Only the former is involved in passing from
oriented to unoriented Euclidean classification.  We define it at the level
of the complete law-valued datum.

Let \(J(x,y)=(x,-y)\) and \(r(u,v)=(-u,-v)\).  For a signed channel-law
collection \(\mathscr D\), write \(\mathscr R_0\mathscr D\) for the collection
obtained by replacing, for every measured edge, both contact types and every
retained offset,
\begin{equation}\label{eq:v28-law-reflection}
 \Pi^{e,\mathrm{end},\infty}_{bb,d}
 \quad\hbox{by}\quad r_*\Pi^{e,\mathrm{end},\infty}_{bb,d}.
\end{equation}
The same involution is used throughout the collection.  The obstacle labels,
the directed edge labels, the marked deck elements and the onset gaps are
unchanged.  In particular this operation does not reverse the directed
channel \((a,b,\ell)\).  For a density it reads
\(f^r_{e,b,d}(u,v)=f_{e,b,d}(-u,-v)\).  When residual time is present its
coordinate is left fixed, and when a full transcript is present its designs,
counts, failure symbols and stopping rule are transported together with the
same change of endpoint coordinates.

The orientation of a marked lattice requires some care.  In the preceding
signed formulation \(\iota\) is required to preserve the orientation of the
marked basis.  A reflection of the physical table does not preserve this
sector.  To describe the common reflection action, let
\(\epsilon\in\{+1,-1\}\) denote the orientation character of \(\iota\).
Write \(\mathscr D_\epsilon=(\epsilon,\mathscr D)\), and let
\(\mathfrak G^\epsilon_{\rm per}(\mathscr D)\) be the gluing space defined by
\eqref{eq:v23-periodic-gluing}--\eqref{eq:v23-periodic-gluing-constraint},
except that the isometry \(\iota\) has orientation character \(\epsilon\).
The placements \(A_e\) are still proper Euclidean motions of the consistently
oriented channel frames.  The original signed statement is the
\(\epsilon=+1\) sector.  The proof of its classification applies in either
sector: it uses the incidence equations and Euclidean isometry of \(\iota\),
not the sign of its determinant.

The character \(\epsilon\) records the common orientation convention relative
to the marked lattice; it is not a supplied lattice angle or a laboratory
registration.  In an unoriented representative it must be transported with
the transverse convention, rather than erased independently.

\begin{definition}[Common-orientation datum]
\label{def:v28-common-orientation-datum}
On the complete marked signed data define
\begin{equation}\label{eq:v28-common-involution}
 \mathscr R(\epsilon,\mathscr D)
   =(-\epsilon,\mathscr R_0\mathscr D),
 \qquad \mathscr R^2=\mathrm{id}.
\end{equation}
The common-orientation datum is the two-element orbit
\[
 \overline{\mathscr D}
   =\{(\epsilon,\mathscr D),
             (-\epsilon,\mathscr R_0\mathscr D)\}.
\]
Thus it retains all relative sign choices, including their relation to the
marked lattice convention, but not a choice of the one common orientation.
It is an orbit of the entire collection of probability laws, not a mixture
of those laws and not a pushforward under pointwise absolute values.
\end{definition}

\begin{lemma}[Equivariance of recovery and gluing]
\label{lem:v28-reflection-equivariance}
Reflection induces a well-defined involution between the two signed gluing
spaces belonging to \(\overline{\mathscr D}\).  On representatives it is
\begin{equation}\label{eq:v28-reflected-gluing}
 \bigl(\iota,(A_e),(C_{e,\sigma})\bigr)
 \longmapsto
 \bigl(J\iota,(JA_eJ),(JC_{e,\sigma})\bigr).
\end{equation}
Reflected boundary images are given their canonical positive boundary
orientation.  The recovered quotient obstacle images are \(JC_r\), and all
labelled deck translations are transported by the same \(J\).
\end{lemma}

\begin{proof}
A reflected contact graph is \(\psi_b^r(y)=\psi_b(-y)\).  Reflection preserves
flight lengths, the reflection law, arclength and the collision flux measure.
Sending every endpoint coordinate to its negative therefore gives
\eqref{eq:v28-law-reflection}; the gap and residual-time coordinate do not
change.  Consequently the signed inverse and analytic globalization recover
exactly the reflected boundary images.  Equivalently the half-line actions
and their graph jets satisfy
\begin{equation}\label{eq:v28-reflected-jets}
 S_b^r(u)=S_b(-u),\qquad
 (q_{b,n})^r=(-1)^nq_{b,n}.
\end{equation}
These identities also follow directly by reflecting each finite stationary
bridge and then taking the half-line limit.  They do not require convergence
of an infinite Taylor series for a smooth graph.

Since \(\det J=-1\), the realization \(J\iota\) lies in the opposite lattice
sector while \(JA_eJ\in\mathrm{SE}(2)\).  The identity
\(J\tau_v=\tau_{Jv}J\) gives, incidence by incidence,
\begin{equation}\label{eq:v28-incidence-equivariance}
 \tau_{-J\iota(\nu)}(JA_eJ)(JC_{e,\sigma})
       =J\bigl(\tau_{-\iota(\nu)}A_eC_{e,\sigma}\bigr).
\end{equation}
Thus every gluing equation is preserved.  Reflection also preserves
disjointness, strict convexity, facing contacts, flight lengths and the
selected channel geometry; hence it preserves admissibility.  Applying the
construction twice gives the original gluing.

Finally a different proper representative \(B\circ A_e\), with lattice
realization \(R_B\iota\), changes the reflected representative by the single
proper motion \(JBJ\).  Formula~\eqref{eq:v28-reflected-gluing} is therefore
well defined after the proper Euclidean quotient.
\end{proof}

\begin{theorem}[Unoriented marked periodic classification]
\label{thm:v28-unoriented-classification}
Under the analytic and measured-network hypotheses of
Theorem~\ref{thm:v23-intrinsic-table-rigidity}, let
\(\overline{\mathscr D}\) be the common-orientation datum of
Definition~\ref{def:v28-common-orientation-datum}.  Define
\begin{equation}\label{eq:v28-quotient-gluing-space}
 \overline{\mathfrak G}_{\rm per}(\overline{\mathscr D})
 =\left(
    \mathfrak G^\epsilon_{\rm per}(\mathscr D)
    \sqcup
    \mathfrak G^{-\epsilon}_{\rm per}(\mathscr R_0\mathscr D)
   \right)\big/\langle\mathscr R\rangle .
\end{equation}
The labelled periodic tables realizing this datum, modulo one common element
of \(\mathrm E(2)\) acting on the lattice and every obstacle image, are in
bijection with
\(\overline{\mathfrak G}_{\rm per}(\overline{\mathscr D})\).
In particular, a singleton signed gluing space gives a singleton unoriented
gluing space and determines the table up to one common Euclidean isometry.
The conclusion applies to a single positive offset at each contact type
whenever the signed reconstruction is made using the single-offset inverse.
\end{theorem}

\begin{proof}
Choose either representative of \(\overline{\mathscr D}\).  The signed
classification, in its indicated lattice sector, identifies proper-motion
classes of realizing tables with the corresponding signed gluing space.
Lemma~\ref{lem:v28-reflection-equivariance} identifies the classifications
for the two representatives by the same reflection of the entire table.

More explicitly, if two realizing tables are related by a proper Euclidean
motion, the signed classification identifies their signed gluing classes.
If they are related by an improper motion, write that motion as \(B J\) with
\(B\in\mathrm{SE}(2)\).  Reflection first interchanges the two signed data
representatives by \(\mathscr R\), and \(B\) then identifies their proper
classes.  Conversely, equality in \eqref{eq:v28-quotient-gluing-space} is
represented either by equality in one signed gluing space or by equality
after its common reflected identification.  The inverse signed
classification yields respectively a proper or an improper global Euclidean
motion of the realizing tables.  This proves both directions of the
bijection.  Reflection maps the two signed gluing spaces bijectively, so the
singleton assertion follows.  The argument uses only the signed inverse,
not the number of offsets used to recover its actions.
\end{proof}

\begin{corollary}[Holonomy and fixed-order stability under the quotient]
\label{cor:v28-quotient-holonomy-stability}
The signature-rigid gluing criteria are invariant under the common
reflection.  If two independent marked cycle columns form \(M\), their
recovered holonomy columns form \(V\), and \(L=VM^{-1}\), then
\begin{equation}\label{eq:v28-lattice-equivariance}
 V^r=JV,\qquad L^r=JL,\qquad (L^r)^{\mathsf T}L^r=L^{\mathsf T}L.
\end{equation}
Thus the common-orientation quotient retains the recovered Gram form and
scale.  Any fixed-order stability estimate for the signed inverse, valid for
the representative pairs used in reflection-invariant orbit metrics,
descends with the same constant to those metrics.
\end{corollary}

\begin{proof}
Canonical positive arclength reverses under reflection.  Hence the curvature
signature transforms as
\( (\kappa,\kappa',\kappa'',\ldots)
\mapsto(\kappa,-\kappa',\kappa'',\ldots)\).
Equality and uniqueness of signature matches are preserved.  Every incidence
congruence is conjugated by \(J\), as is its product around a cycle.  Since
\(J\tau_vJ=\tau_{Jv}\), the holonomy vector becomes \(Jv\); marked cycle
columns are unchanged.  Formula~\eqref{eq:v28-lattice-equivariance} follows
from \(J^{\mathsf T}J=I\).  A rooted rigid spanning tree is carried to the
same rooted rigid combinatorial tree with all its matches reflected.

For the stability assertion, let \(F\) be the signed inverse between the
relevant fixed-order data and reconstruction spaces, with reflection-invariant
metrics \(d\) and \(d'\), and suppose
\(d'(F(D),F(E))\le C_Md(D,E)\) for every representative pair entering the orbit comparison.
If the estimate is formulated sectorwise, use only same-sector
representative comparisons; their reflected estimates have the same constant.
Define the orbit distances by
\[
 \bar d([D],[E])=\min_{\eta\in\{0,1\}}d(D,\mathscr R^\eta E),
 \qquad
 \bar d'([X],[Y])=\min_{\eta\in\{0,1\}}d'(X,\mathscr R^\eta Y).
\]
For a sectorwise metric the minima are restricted to the nonempty set of
same-sector comparisons.  Choose a minimizing \(\eta\) and use equivariance
of \(F\).  Then
\[
 \bar d'([F(D)],[F(E)])
 \le d'(F(D),F(\mathscr R^\eta E))
 \le C_M\bar d([D],[E]).
\]
For contact jets and densities on symmetric compact boxes, reflection
invariance follows from \eqref{eq:v28-reflected-jets} and the usual
fixed-order supremum norms.  This statement retains the order-dependent
constants and hypotheses of the signed estimates; it asserts no
order-independent stability for analytic continuation.
\end{proof}

\begin{remark}[Law-valued orientation orbits are not folded records]
\label{rem:v28-folded-records}
For a continuously distributed endpoint pair, the pushforward under
\((u,v)\mapsto(|u|,|v|)\) has density, on the open positive quadrant,
\begin{equation}\label{eq:v28-folded-density}
 f_{\rm abs}(a,b)=f(a,b)+f(-a,b)+f(a,-b)+f(-a,-b).
\end{equation}
The coordinate axes have zero mass in this setting.  This sum is not the
law-valued orbit \(\{f,f\circ r\}\).  Even pushing each observed pair to its
orbit under simultaneous reversal is a sample-space coarsening, not an orbit
of the complete probability law.  Randomly symmetrizing a law is a further
operation, and independently reversing different channel collections gives
an action of \(\{\pm1\}^{E}\), rather than the single diagonal group used
above.

For example, on \([-1,1]^2\), the positive normalized symmetric densities
\[
 f_1(u,v)=\tfrac14\bigl(1+\alpha(u+v)\bigr),\qquad
 f_2(u,v)=\tfrac14\bigl(1+2\alpha(u+v)\bigr),
 \qquad 0<\alpha<\tfrac14,
\]
have the same folded density, identically one on \((0,1)^2\), but their
law-valued common-reversal orbits are different: the coefficient of \(u+v\)
can only change sign under \(r\), not double.  This example distinguishes
the observation maps in a class of densities.  It is not a claim that these
two densities are realized by physical billiard tables, nor a physical
nonidentifiability theorem for folded billiard observations.  All rigidity
statements above use the specified law-valued common-orientation datum.
\end{remark}
